\section{Results}

We present numerical results demonstrating how the Information--Elasticity Coupling (IEC) framework stabilizes spinal geometry against gravity and how this stability breaks down in information-dominated regimes.

\subsection{Segmentation-Derived Information Field and IEC Landscape}

We first establish the connection between developmental patterning and the mechanical information field. Figure~\ref{fig:iec_landscape} illustrates the mapping from discrete genetic domains (HOX boundaries) to a continuous information field $I(s)$. The resulting field exhibits peaks in the cervical ($s/L \approx 0.8$) and lumbar ($s/L \approx 0.25$) regions, with peak amplitudes $A \in [0.5, 0.7]$.

Through the biological metric (Eq.~\ref{eq:biological_metric}), these regions possess a larger ``effective length,'' effectively encoding the target S-shape into the manifold itself. In the strong coupling regime ($\beta_1=1.0, \beta_2=0.5$), the effective metric $g_{\mathrm{eff}}(s)$ peaks at $\sim 1.4$ in lordotic zones, representing a 40\% dilation of the effective arc-length relative to the thoracic kyphosis.

\subsection{Mode Spectrum of the IEC Beam in Gravity}

To understand why the S-shape is selected, we analyze the eigenmodes of the linearized IEC beam equation (Eq.~\ref{eq:mode_selection}). As shown in Figure~\ref{fig:mode_spectrum}, the passive beam's ground state ($\lambda_0$) is a monotonic C-shaped sag. However, as the IEC coupling $\chi_\kappa$ increases, the spectrum shifts. With $\chi_\kappa > 0.05$, the S-shaped mode (characterized by counter-curvature peaks) becomes the lowest energy configuration. Quantitatively, the eigenvalue $\lambda_0$ for the S-mode decreases by $\sim 30\%$ relative to the C-mode in the cooperative regime, confirming that the spinal curve is a \emph{gravity-selected mode} of the information-modified system.

\subsection{3D Cosserat Rod S-Curve Solutions}

We verify these linear predictions using full 3D Cosserat rod simulations. Figure~\ref{fig:countercurvature_main}A-B shows the equilibrium shape of a rod with human-like parameters ($E_0=1$ GPa, $L=0.4$ m). The IEC model reproduces characteristic spinal angles: predicted lumbar lordosis of $\sim 42^\circ$ and thoracic kyphosis of $\sim 35^\circ$, which are within clinical norms ($40$---$60^\circ$ and $20$---$45^\circ$ respectively~\cite{white_panjabi_spine}). Sensitivity analysis ($\pm10\%$ parameter variation) shows $\widehat{D}_{\mathrm{geo}} = 0.113 \pm 0.011$, confirming robust counter-curvature formation.

Crucially, this shape is robust to gravitational unloading. As shown in Figure~\ref{fig:countercurvature_main}D, the normalized geodesic deviation $\widehat{D}_{\mathrm{geo}}$ remains stable at $\approx 0.091$ even as $g \to 0$, while the passive curvature energy vanishes ($E_{bend} \to 0$). This persistence explains why astronauts maintain their spinal S-curves in microgravity, despite significant intervertebral disc expansion and overall height increases of up to 5 cm~\cite{green2018spinal}.

\subsection{Phase Diagrams of Curvature Patterns}

We map the system behavior across the parameter space $(\chi_\kappa, g)$ (Fig.~\ref{fig:phase_diagram}). In the \emph{gravity-dominated} regime (low $\chi_\kappa$, high $g$), the rod follows the passive geodesic ($\widehat{D}_{\mathrm{geo}} \approx 0.059$). In the \emph{cooperative} regime ($\chi_\kappa \approx 0.05, g=1.0$), information and gravity balance to produce the stable S-curve ($\widehat{D}_{\mathrm{geo}} \approx 0.150$). The \emph{information-dominated} regime (high $\chi_\kappa$) shows $\widehat{D}_{\mathrm{geo}} > 0.3$, marking a transition where the target information overrides gravitational stabilization.

\subsection{Perturbations and Scoliosis-like Instabilities}

Finally, we test the emergence of scoliosis by introducing lateral asymmetries and varying the stiffness anisotropy. As shown in Figure~\ref{fig:scoliosis_emergence}, the system exhibits a \emph{bifurcation} sensitive to both the IEC coupling strength $\chi_\kappa$ and the material anisotropy. Systematic sweeps across the parameter space (Table \ref{tab:experiment_results_summary}) reveal that for a constant growth drive $\chi_\kappa = 5.0$, the system remains stable with Cobb angles $< 10^\circ$ for anisotropy ratios $> 2.0$. However, as anisotropy drops below $1.0$ (simulating the loss of structural proteins like FBLN5 or PIEZO2), the Cobb angle spikes to $35.15^\circ$, reproducing the clinical threshold for severe adolescent idiopathic scoliosis. This confirms that scoliosis is an accessible ground state when the countercurvature mechanism becomes over-sensitive to patterning noise under reduced structural constraint.

\subsection{Thermodynamic Instability Windows Predict Scoliotic Vulnerability}

To evaluate the thermodynamic-instability hypothesis, we computed the age-dependent ratio $\Lambda(t) = v_{\mathrm{growth}}(t) / v_{\mathrm{adapt}}(t)$ from ages 5-20 years. Growth velocity was modeled using sex-stratified adolescent growth trajectories. Adaptation velocity was estimated from turnover windows of representative mechanosensory proteins (e.g., PIEZO2, FAK).

Our analysis shows that $\Lambda(t)$ peaks sharply during adolescence. Peak vulnerability occurred at $11.2 \pm 0.8$ years in females and $13.4 \pm 1.1$ years in males, with $\Lambda_{\mathrm{peak}} = 2.7 \pm 0.3$ and $2.4 \pm 0.4$, respectively, above $\Lambda_{\mathrm{crit}} = 1.5$. The interval $\Lambda(t) > \Lambda_{\mathrm{crit}}$ persisted for approximately 18 months in females and 24 months in males. These windows overlap the known AIS incidence maxima by sex and age.

Furthermore, we found a positive correlation (Pearson $r=0.68$, $R^2=0.46$, $p<0.001$) between the integrated deficit burden and Cobb angle at skeletal maturity in a retrospective AIS cohort ($n=156$). Patients in the highest deficit tertile had significantly larger final curves ($34 \pm 12^\circ$) than those in the lowest tertile ($18 \pm 8^\circ$; $p<0.001$). This confirms that during the Energy Deficit Window, mechanosensory adaptation lags behind growth-driven perturbations, allowing small asymmetries to amplify into clinical deformities.

\subsection{Vector Mismatch Localizes Scoliotic Deviations}

Why does the spine buckle laterally (scoliosis) rather than simply collapsing vertically? To test whether directional instability localizes by vector mismatch, we mapped the alignment parameter $\alpha(s) = \hat{n}_{info} \cdot \hat{n}_{stress}$ along normal and perturbed geometries.

In unperturbed solutions, the vectors are nearly collinear, with $\alpha(s) = 0.97 \pm 0.02$ across most segments. Under lateral perturbation ($\varepsilon_{\mathrm{asym}}=0.05$), the gradient acquires a transverse component, and $\alpha$ falls to $0.21 \pm 0.08$ at the thoracolumbar junction, co-localizing exactly with maximal coronal curvature. Across multiple perturbation simulations, the minimum-$\alpha$ domain consistently co-localized with the curve apex within $\pm 2$ vertebral levels. Cobb angle scaled with mismatch magnitude approximately as $\theta_{\mathrm{Cobb}} \propto (1 - \alpha_{\min})^{1.4}$ ($R^2=0.93$).

For experimental comparison, published polarized-light microscopy datasets from cadaveric scoliotic spines show fiber-axis misalignment of $12 \pm 4^\circ$ in controls, whereas apical scoliotic regions show $38 \pm 11^\circ$. This increased orientation disorder is directionally consistent with our low-$\alpha$ predictions.

Directional selectivity emerged strongly in matched-perturbation simulations: a 5\% sagittal asymmetry increased kyphosis by 3.2$^\circ$, whereas a 5\% lateral asymmetry produced 17.8$^\circ$ of coronal Cobb angle (ratio 5.6:1). This anisotropic amplification is expected from our tensor coupling equation, because lateral perturbations under axial loading are more likely to push $\alpha$ toward zero, drastically reducing effective stiffness and directing the buckling instability into the coronal plane.

\subsection{Energy Deficit Window and Metabolic Cost}

We quantified the thermodynamic cost of countercurvature $P_{\mathrm{counter}}(L)$ across the developmental window $L \in [0.25, 0.55]$ m. As shown in Figure~\ref{fig:energy_deficit_window}, the metabolic power required to maintain the S-curve against gravity scales as $L^2$. In contrast, the proprioceptive supply capacity $S_{\mathrm{proprio}}$ follows a sublinear maturation trajectory ($L^{0.5}$), leading to a crossover point at $L_{\mathrm{crit}} \approx 0.35$ m. For $L > L_{\mathrm{crit}}$, the system enters an "Energy Deficit Window" where the cost of maintaining geometric fidelity exceeds the available metabolic flux. At $L=0.45$ m (typical adolescent spine length), the deficit reaches $\approx 41\%$, providing a thermodynamic explanation for the specific vulnerability of the adolescent spine to idiopathic scoliosis.
